
We measured the clusterability of spurious features in SSL embeddings trained on the Waterbirds spurious correlation dataset. Under the tested conditions (SimCLR and LA-SSL training, ResNet-50 architecture, 20-epoch proof-of-concept experiments), observed AMI values were approximately 0.28, below the 0.4 threshold used to indicate reliable clusterability. Linear separability remained high (AUC $\approx$ 0.98) despite low clusterability, indicating that these are dissociable properties. LA-SSL increased AMI by 2.04\% rather than decreasing it by the hypothesized 30\%, while preserving linear separability.

These results do not support the hypothesis that standard SSL creates discrete, geometrically separable spurious clusters within 20 epochs of training. The mechanisms tested (M1: InfoNCE cluster formation, M2: clusterability predicts efficacy, M3: LA-SSL disperses clusters) were not supported under the tested conditions. Extended training experiments (100 epochs) are needed to determine whether these findings hold at scale or whether cluster structure emerges with longer training duration.

The primary contributions are: (1) empirical measurements of clusterability (AMI) in SSL embeddings on a spurious correlation dataset, (2) demonstration that linear separability and discrete clusterability are dissociable under the tested conditions, and (3) validated implementations of cluster-based mechanisms (43/43 tests passing for h-e1, 5/5 tests passing for h-m-integrated, 100\% SDD compliance for both) suitable for future experimental work.
